%
% Halaman Abstract
%
% @author  Andreas Febrian @version 2.2.0 @edit by Ichlasul Affan
%

\chapter*{ABSTRACT}
\singlespacing

\noindent \begin{tabular}{l l p{11.0cm}} \ifx\blank\npmDua Name&: & \penulisSatu
	\\
		Study Program&: & \studyProgramSatu \\
	\else
		Writer 1 / Study Program&: & \penulisSatu~/ \studyProgramSatu\\
		Writer 2 / Study Program&: & \penulisDua~/ \studyProgramDua\\
	\fi
	\ifx\blank\npmTiga\else Writer 3 / Study Program&: & \penulisTiga~/
		\studyProgramTiga\\
	\fi
	Title&: & \judulInggris \\
	Counselor&: & \pembimbingSatu \\
	\ifx\blank\pembimbingDua
	\else
		\ &\ & \pembimbingDua \\
	\fi
	\ifx\blank\pembimbingTiga
	\else
		\ &\ & \pembimbingTiga \\
	\fi
\end{tabular} \\

\vspace*{0.5cm}

\noindent This internship was carried out at PT TASPEN (Persero) with the scope
of work being the development of the Easy Mobile internal application, which
supports employee activities such as submitting leave requests and work orders
(SPT). During the internship, the intern served as a Mobile Developer Intern in
the Information Technology Division with the responsibility of developing the
SPT Request and Leave Request features using Flutter, as well as utilizing
Figma, GitLab, Postman, and Android Studio as supporting tools. The
implementation process faced challenges in the form of limited API documentation
and complex user flows, but these were resolved through intensive communication
with mentors, independent analysis, and the adoption of clean architecture
practices. This activity not only strengthened technical skills in mobile
application development but also honed non-technical skills such as
communication, teamwork, adaptability, and time management, which are essential
assets in the professional workplace. \\

\vspace*{0.2cm}

\noindent Key words: \\ mobile development, Flutter, front-end, enterprise
application \\

\setstretch{1.4}
\newpage
